% Example output — Dr. Jordan Chen applying to Whitfield University
% Generated by claude-resume-kit for demonstration purposes
% This is a fictional researcher; all data is fabricated.
% Format: CV (5-page, cv.cls) — matches example_session_file.md's documented output.

\documentclass{cv}
\usepackage{hyperref}
\usepackage{fontawesome}
\usepackage{enumitem}
\usepackage{tikz}
\usepackage{graphicx}
\hypersetup{
    colorlinks = true,
    linkcolor = [rgb]{0.1,0.1,0.6},
    citecolor = [rgb]{0.4,0.4,0.4},
    urlcolor = [rgb]{0.0,0.0,0.7},
    pdftitle = {Jordan Chen, Ph.D. - Academic CV},
    pdfauthor = {Jordan Chen, Ph.D.}
}
\usepackage{xcolor}
\usepackage[version=4,arrows=pgf-filled]{mhchem}
\usepackage[includefoot,left=0.75in,top=0.76in,right=0.75in,bottom=0.7in,textwidth=7.5in,textheight=10.9in]{geometry}
\usepackage{fancyhdr}
\pagestyle{fancy}
\fancyhf{}
\renewcommand{\headrulewidth}{0pt}
\rfoot{\thepage/\pageref{LastPage}}

\name{Jordan Chen, Ph.D.}
\address{\faEnvelope~\href{mailto:jordan.chen@email.com}{jordan.chen@email.com}}
\address{\faMapMarker~Richland, WA $|$ \faLinkedin~\href{https://linkedin.com/in/jordanchen}{LinkedIn} $|$ \includegraphics[height=1em]{orcid.png}~\href{https://orcid.org/0000-0002-1234-5678}{ORCID: 0000-0002-1234-5678} $|$ \includegraphics[height=1em]{GS.png}~\href{https://scholar.google.com/citations?user=XXXXXXXXX}{Google Scholar}}

\begin{document}
\vspace{-0.03cm}

\begin{rSection}{Research Summary}
Computational biologist with 8+ years spanning molecular dynamics simulation and ML-guided protein engineering, bridging biomolecular simulation and data-driven design for enzyme engineering and drug discovery. Fine-tuned the ESM-2 protein language model on 45K experimental stability measurements to screen 8,500 enzyme variants at 3,000$\times$ experimental throughput, with 5 hits confirmed by collaborators via differential scanning calorimetry. Co-developed an open-source transfer learning framework adopted by 4 external research groups, and built enhanced-sampling protocols for protein folding and ligand binding that now underpin the group's simulation infrastructure. Author of 15 publications (7 first-author, 280+ citations) in venues including ACS Catalysis, J. Chem. Theory Comput., and J. Med. Chem.
\end{rSection}
\vspace{-0.01cm}

\begin{rSection}{Education}
\textbf{Ph.D., Biomedical Engineering} \hfill {Aug 2018 -- Jul 2023}\\
Westfield Institute of Technology, Westfield, MA \hfill GPA: 3.92/4.00\\
\textit{Dissertation:} Enhanced Sampling Methods for Protein Folding and Ligand Binding Free Energy Prediction\\
\textit{Advisor:} Prof. L. Alvarez

\textbf{B.S., Biochemistry (Honors)} \hfill {Aug 2014 -- May 2018}\\
Eastgate University, Portland, OR \hfill GPA: 3.87/4.00
\end{rSection}
\vspace{-0.0cm}

\begin{rSection}{Technical Expertise}

\begin{skillgroup}{Molecular Simulation \& Modeling}
\skilldash{\textbf{GROMACS}, OpenMM, AMBER -- metadynamics, replica exchange MD, free energy perturbation}
\skilldash{AlphaFold2, Rosetta, AutoDock Vina -- structure prediction, homology modeling, docking}
\skilldash{CHARMM36m, AMBER ff19SB, OPLS-AA/M -- force field benchmarking for disordered proteins}
\end{skillgroup}

\begin{skillgroup}{Machine Learning \& Data Science}
\skilldash{\textbf{Protein language models} (ESM-2, 650M), graph neural networks, transfer learning}
\skilldash{PyTorch, scikit-learn, BioPython -- fine-tuning, embedding extraction, feature engineering}
\skilldash{Regression, classification, cross-validation, Spearman/RMSE benchmarking, dataset curation}
\end{skillgroup}

\begin{skillgroup}{Analysis \& Visualization}
\skilldash{MDAnalysis, ProDy, PyMOL, matplotlib, seaborn -- trajectory analysis, structural visualization}
\skilldash{PostgreSQL, pandas, NumPy -- curated stability databases with automated quality filters}
\end{skillgroup}

\begin{skillgroup}{Domain Expertise}
\skilldash{Protein engineering, enzyme thermostability, folding thermodynamics, drug discovery}
\skilldash{Intrinsically disordered proteins, ligand binding free energy, biocatalysis, directed evolution}
\end{skillgroup}

\begin{skillgroup}{Programming \& HPC}
\skilldash{Python, Bash, SQL -- scientific computing, data pipelines, automated analysis workflows}
\skilldash{SLURM, Snakemake, Git, DVC -- HPC job scheduling, workflow automation, reproducible research}
\end{skillgroup}

\end{rSection}
\vspace{-0.0cm}

\begin{rSection}{Research Experience}

\begin{rSubsection}{Postdoctoral Research Associate}{\textcolor{black!60}{Aug 2023 -- Present}}{Lakewood University, Department of Bioengineering}{PI: Prof. K. Holmberg}

\subtheme{ML-Guided Enzyme Stability Screening}
\item Fine-tuned the ESM-2 protein language model on 45K experimental melting temperatures, achieving 0.82 Spearman correlation and enabling 3,000$\times$ throughput screening of 8,500 enzyme variants.
\item Co-developed a transfer learning framework reducing labeled training data requirements by 60\% across five enzyme families, released as an open-source tool with 200+ GitHub stars.
\subtheme{Enzyme Solvent Tolerance for Green Chemistry}
\item Extended the protein language model to predict enzyme solvent tolerance across eight organic co-solvent systems, validating against 50-ns explicit-solvent MD for 80 enzyme variants.
\subtheme{Automated Simulation Infrastructure}
\item Automated the sequence-to-simulation pipeline using the Snakemake workflow manager, reducing per-variant setup time from 4 hours to 10 minutes for six researchers.

\end{rSubsection}

\begin{rSubsection}{Ph.D. Researcher}{\textcolor{black!60}{Aug 2018 -- Jul 2023}}{Westfield Institute of Technology}{Advisor: Prof. L. Alvarez}

\subtheme{Enhanced Sampling for Protein Folding}
\item Developed a metadynamics-based enhanced sampling protocol for protein folding free energy landscapes, predicting folding temperatures within 8 K of experiment across six proteins.
\subtheme{Ligand Binding Free Energy Calculations}
\item Calculated relative binding free energies for 40 congeneric ligand pairs via free energy perturbation, achieving 0.9 kcal/mol RMSE against experimental IC50 data.
\subtheme{Thermostability Database Infrastructure}
\item Built a curated protein thermostability database integrating 12,000 experimental melting temperatures from 3 public sources, with quality filters adopted by 8 lab members.

\end{rSubsection}

\begin{rSubsection}{Undergraduate Research Assistant}{\textcolor{black!60}{May 2016 -- Jul 2018}}{Eastgate University}{}

\item Performed homology modeling and 100-ns MD simulations of 4 mutant lysozyme variants, identifying destabilizing cavity mutations consistent with published experimental data.
\item Built Python analysis scripts for automated hydrogen bond occupancy tracking across 500-ns aggregate trajectories, adopted by 3 lab members for stability projects.

\end{rSubsection}

\end{rSection}

\begin{rSection2}{Fellowships \& Honors}
\item \textbf{NSF Graduate Research Fellowship}, National Science Foundation (2019). Three-year fellowship supporting doctoral research in computational protein engineering.
\item \textbf{Best Oral Presentation}, Westfield Biophysics Symposium (2022). Enhanced sampling methods for protein folding thermodynamics.
\item \textbf{Dean's Teaching Award}, Westfield Institute of Technology (2021). Outstanding TA in computational biology.
\item \textbf{Outstanding Poster Award}, Gordon Research Conference on Proteins (2023). ML-guided enzyme thermostability screening.
\end{rSection2}

\begin{rSection}{Publications}
\textit{15 peer-reviewed articles $|$ 280+ citations $|$ h-index: 9} \hfill \href{https://scholar.google.com/citations?user=XXXXXXXXX}{Google Scholar}\\
\vspace{0.07cm}
Selected publications shown below; full list available on Google Scholar.
\vspace{0.15cm}

\textbf{Published}
\begin{enumerate}[leftmargin=1.5em, labelsep=0.5em, itemsep=0.1em]
\item[6.] \textbf{J. Chen}, R. Nakamura, S. Patel, K. Holmberg, M. Rivera. ``Deep Learning-Guided Screening of Thermostable Enzyme Variants for Industrial Biocatalysis.'' \textit{ACS Catalysis} (2025).
\item[5.] \textbf{J. Chen}, M. Rivera, K. Holmberg. ``Transfer Learning from Protein Language Models for Low-Data Enzyme Property Prediction.'' \textit{Bioinformatics} (2024).
\item[4.] \textbf{J. Chen}, L. Alvarez. ``Ligand Binding Free Energies via Enhanced-Sampling FEP for Three Drug Target Families.'' \textit{J. Med. Chem.} (2023).
\item[3.] \textbf{J. Chen}, L. Alvarez. ``Metadynamics Protocol for Protein Folding Free Energy Landscapes.'' \textit{J. Chem. Theory Comput.} (2022).
\item[2.] \textbf{J. Chen}, P. Kowalski, L. Alvarez. ``Force Field Benchmarking for Intrinsically Disordered Protein Ensembles.'' \textit{J. Chem. Theory Comput.} (2021).
\item[1.] \textbf{J. Chen}, P. Kowalski, L. Alvarez. ``Curated Thermostability Database for ML-Ready Protein Engineering Benchmarks.'' \textit{Bioinformatics} (2021).
\end{enumerate}

\textbf{Under Review}
\begin{enumerate}[leftmargin=1.5em, labelsep=0.5em, itemsep=0.1em]
\item[--] \textbf{J. Chen}, T. Yamamoto, K. Holmberg. ``Predicting Enzyme Solvent Tolerance with Fine-Tuned Protein Language Models.'' \textit{Proteins: Struct., Funct., Bioinf.} Under review.
\end{enumerate}
\end{rSection}

\begin{rSection2}{Selected Presentations}
\item \textbf{J. Chen}. ``ML-Guided Screening of Thermostable Enzyme Variants.'' \textit{Gordon Research Conference on Proteins}, Ventura, CA (2023). Poster.
\item \textbf{J. Chen}, M. Rivera, K. Holmberg. ``Transfer Learning from Protein Language Models for Low-Data Enzyme Property Prediction.'' \textit{ACS National Meeting}, San Francisco, CA (2024). Oral.
\item \textbf{J. Chen}, L. Alvarez. ``Enhanced Sampling Protocols for Protein Folding Free Energy Landscapes.'' \textit{Biophysical Society Annual Meeting}, San Diego, CA (2022). Oral.
\item \textbf{J. Chen}, P. Kowalski, L. Alvarez. ``Force Field Benchmarking for Intrinsically Disordered Protein Ensembles.'' \textit{AIChE Annual Meeting}, Boston, MA (2021). Poster.
\item \textbf{J. Chen}. ``Automated Screening Pipelines for ML-Guided Enzyme Engineering.'' \textit{Lakewood University Computational Biology Seminar Series} (2024). Invited talk.
\end{rSection2}

\begin{rSection2}{Mentorship \& Professional Service}
\item \textbf{Graduate mentorship:} Mentored 3 graduate students on protein ML pipelines and MD simulation workflows; one student co-authored a peer-reviewed publication within 8 months of joining.
\item \textbf{Teaching:} Recognized with the Dean's Teaching Award (2021) for outstanding TA performance in computational biology coursework.
\item \textbf{Lab training adoption:} Analysis scripts and quality-filter protocols adopted by 3--8 lab members across three research groups for ongoing protein stability projects.
\end{rSection2}

\end{document}
